\section{Билет 5. Унимодальные функции. Одномерная минимизация без производной}

\ticketq{Унимодальные функции. Методы одномерной минимизации, не использующие
вычисление производной функции цели.}

\paragraph{Постановка.}
$f(x) \to \min$, $x \in [a,b]$. Методы \textbf{нулевого порядка} используют
только значения $f(x_1), f(x_2), \dots$ и не используют $f'(x)$. Применяются
самостоятельно либо как подбор шага $\alpha_k$ в многомерных методах:
$\min_\alpha f(x_k + \alpha d_k)$.

\paragraph{Унимодальность.}
$f$ \textbf{унимодальна} на $[a,b]$, если существует единственная точка
минимума $x_*$, причём $f$ убывает на $[a, x_*]$ и возрастает на $[x_*, b]$:
\[
  x_1 < x_2 \le x_* \Rightarrow f(x_1) \ge f(x_2) ,
  \qquad
  x_* \le x_1 < x_2 \Rightarrow f(x_1) \le f(x_2) .
\]
\textbf{Интервал неопределённости} — отрезок $[a_k, b_k] \ni x_*$,
$L_k = b_k - a_k$. Останов при $L_k \le \delta$, тогда
$x_* \approx (a_k + b_k)/2$.

\paragraph{Лемма (об исключении части отрезка).}
\textit{Пусть $f$ унимодальна на $[a,b]$ и $x_1 < x_2$. Тогда}
\[
  f(x_1) \ge f(x_2) \ \Rightarrow\ x_* \notin [a, x_1) ,
  \qquad
  f(x_1) \le f(x_2) \ \Rightarrow\ x_* \notin (x_2, b] .
\]

\textsc{Док-во.} Первое, от противного: $x_* < x_1 < x_2$ $\Rightarrow$ обе
точки правее минимума $\Rightarrow$ по возрастанию справа
$f(x_1) \le f(x_2)$, а по условию $f(x_1) \ge f(x_2)$ $\Rightarrow$
$f(x_1) = f(x_2)$ — невозможно при единственном минимуме на строго
возрастающем участке. Второе симметрично: $x_1 < x_2 < x_*$ $\Rightarrow$
$f(x_1) \ge f(x_2)$ по убыванию слева, что противоречит
$f(x_1) \le f(x_2)$. $\blacksquare$

Вся дальнейшая работа — применение леммы: по двум сравнённым значениям
отбрасывается часть отрезка.

\paragraph{Равномерный поиск.}
$h = (b-a)/n$, $x_i = a + ih$, $i = \overline{0,n}$. Найти $j$ с
$f(x_j) = \min_i f(x_i)$ $\Rightarrow$ $x_* \in [x_{j-1}, x_{j+1}]$, длина
$2h$. Стоит $n+1$ вычислений.

\paragraph{Локализация экстремума.}
То же, но без просмотра всей сетки: идём слева и останавливаемся на первом
$i$ с
\[
  f\big(a+(i-1)h\big) > f\big(a+ih\big) \le f\big(a+(i+1)h\big)
  \ \Rightarrow\ x_* \in \big[a+(i-1)h,\ a+(i+1)h\big] .
\]
Результат тот же, вычислений в среднем меньше: обход прекращается сразу
после минимума.

\textbf{Вариант с удвоением шага (метод Свенна)} — применяется, когда отрезок
заранее не задан. Выбираем $x_0$, $h > 0$; по значениям
$f(x_0-h), f(x_0), f(x_0+h)$ определяем направление убывания и идём в нём с
удвоением шага:
\[
  x_{k+1} = x_k + 2^k h .
\]
После первого возрастания минимум локализован: $x_* \in [x_{k-1}, x_{k+1}]$.
Выигрыш в том, что отрезок находится за $O(\log)$ шагов, а не перебором сетки.

\paragraph{Дихотомия.}
$x_{1,2} = \dfrac{a_k+b_k}{2} \mp \eps$, $\eps > 0$ мало. По лемме:
\[
  f(x_1) \ge f(x_2) \Rightarrow [a_{k+1},b_{k+1}] = [x_1, b_k] ,
  \qquad
  \text{иначе } [a_k, x_2] .
\]
$L_{k+1} = L_k/2 + \eps$ за \emph{два} вычисления $\Rightarrow$ коэффициент
$\sqrt{1/2} \approx 0{,}707$ на одно вычисление. Параметр $\eps$ берётся на
пределе различимости значений $f$.

\paragraph{Бисекция (деление отрезка пополам).}
Три точки $x_1 = a_k + L_k/4$, $x_2 = (a_k+b_k)/2$, $x_3 = a_k + 3L_k/4$:
\[
  f(x_1) < f(x_2) \Rightarrow [a_k, x_2] ; \quad
  \text{иначе } f(x_2) > f(x_3) \Rightarrow [x_2, b_k] ; \quad
  f(x_2) \le f(x_3) \Rightarrow [x_1, x_3] .
\]
Отрезок ровно вдвое короче, а серединой нового становится одна из уже
вычисленных точек $\Rightarrow$ два новых вычисления на итерацию, тот же
коэффициент $0{,}707$.

\paragraph{Золотое сечение.}
$\lambda_k = a_k + \alpha L_k$, $\mu_k = b_k - \alpha L_k$; правило то же:
$f(\lambda_k) \ge f(\mu_k) \Rightarrow [\lambda_k, b_k]$, иначе
$[a_k, \mu_k]$.

\textsc{Вывод $\alpha$.} После сокращения $L_{k+1} = (1-\alpha)L_k$. Чтобы
одна из старых точек стала внутренней точкой нового отрезка (и её значение
переиспользовалось), нужно
\[
  \frac{\alpha}{1-\alpha} = 1 - \alpha
  \ \Rightarrow\ \alpha = (1-\alpha)^2
  \ \Rightarrow\ \alpha^2 - 3\alpha + 1 = 0
  \ \Rightarrow\ \alpha = \frac{3 - \sqrt5}{2} \approx 0{,}382
\]
(второй корень $>1$ не годится). Итог: \emph{одно} новое вычисление $f$ на
итерацию, коэффициент $1-\alpha \approx 0{,}618$.

\paragraph{Метод Фибоначчи.}
$F_0 = F_1 = 1$, $F_{i+1} = F_i + F_{i-1}$; число шагов $n$ фиксируется
заранее и выбирается из требуемой точности: $F_{n} \ge (b-a)/\delta$. На
$k$-м шаге
\[
  \lambda_k = a_k + \frac{F_{n-k-1}}{F_{n-k+1}}L_k ,
  \qquad
  \mu_k = a_k + \frac{F_{n-k}}{F_{n-k+1}}L_k ,
\]
правило сокращения прежнее, и так же переиспользуется одна точка. Итоговая
длина интервала $(b-a)/F_n$ — \textbf{наилучшая возможная} при заданном числе
вычислений $f$, то есть метод оптимален. Золотое сечение — его предельный
вариант: $F_{n-1}/F_n \to 0{,}618$, но ему не нужно знать $n$ заранее.

\paragraph{Сводка (коэффициент сокращения на одно вычисление $f$).}
равномерный поиск $\approx 0{,}794$; дихотомия и бисекция $\approx 0{,}707$;
золотое сечение $\approx 0{,}618$; Фибоначчи — оптимум.

\seefull{гл.~2, \S9--\S10 (численный пример для всех шести методов там же)}
